\chapter{Probability}
\label{ch:probability}

\begin{tcolorbox}[feynbox={Sky}{BBg},title={Before and After}]
\textbf{Before this chapter you need:} fractions and percentages
(\S\ref{sec:percentages}), functions and their range
(\S\ref{sec:functions}--\S\ref{sec:domain}), the exponential and logarithm
(Chapter~\ref{ch:explog}), integrals as accumulated area
(\S\ref{sec:integrals}), vectors and the dot product (\S\ref{sec:vectors}).\\
\textbf{After it you will understand:} what a random variable is, how densities
and probabilities relate, what expectation and variance actually average, why
independence is the assumption market returns violate, and --- the point the whole
thesis turns on --- why the Central Limit Theorem does not rescue you when the
fourth moment does not exist.
\end{tcolorbox}

Probability is the mathematics of uncertainty. This chapter builds it from
counting outcomes up to the two limit theorems, and then shows precisely where
those theorems stop applying to financial returns.

%% ─────────────────────────────────────────────────────────────────────────────
\section{Sets and Events}
\label{sec:sets}

\situation{
Roll one die. Six things can happen. You might care about ``a six'', or about
``an even number'', or about ``more than four''. Each of these is a collection of
outcomes: $\{6\}$, $\{2,4,6\}$, $\{5,6\}$. The question you ask picks out a
collection.
}

\problem{
Before anything can be assigned a probability, we need to say clearly what the
things being assigned probabilities \emph{are}. ``The chance it rains'' is
ambiguous until ``it rains'' is pinned to a definite collection of outcomes, and
ambiguity here produces paradoxes later.
}

\workedarith{
Rolling one die, the complete list of outcomes is
\[
  \Omega = \{1,2,3,4,5,6\} .
\]
Let $A$ = ``even'' $= \{2,4,6\}$ and $B$ = ``more than four'' $= \{5,6\}$.

\medskip
\textbf{Both happen} (intersection): $A \cap B = \{6\}$ --- one outcome.

\textbf{Either happens} (union): $A \cup B = \{2,4,5,6\}$ --- four outcomes.

\textbf{$A$ does not happen} (complement): $A^{c} = \{1,3,5\}$ --- three outcomes.

\medskip
\textbf{Counting check.} $|A| = 3$, $|B| = 2$, $|A \cap B| = 1$, and
\[
  |A \cup B| = |A| + |B| - |A \cap B| = 3 + 2 - 1 = 4 . \quad\checkmark
\]
Subtracting the intersection is necessary because the outcome 6 lies in both sets
and would otherwise be counted twice.
}

\formaldef{
The \textbf{sample space} $\Omega$ is the set of all possible outcomes. An
\textbf{event} is a subset of $\Omega$.

For events $A, B \subseteq \Omega$:
\begin{itemize}
  \item $A \cap B$ (intersection): both occur;
  \item $A \cup B$ (union): at least one occurs;
  \item $A^{c}$ (complement): $A$ does not occur;
  \item $A \cap B = \emptyset$: the events are \textbf{disjoint} or
        \textbf{mutually exclusive} --- they cannot both occur.
\end{itemize}
}

\analogybreak{
A die has six outcomes and every subset is a sensible event. For continuous
quantities this breaks down: if $\Omega$ is the real line, some subsets are so
pathological that no consistent probability can be assigned to them. Measure
theory exists to exclude those, restricting attention to a well-behaved family of
events. Nothing in this thesis meets such a set, but the restriction is why the
formal definition below speaks of a chosen family rather than all subsets.
}

\whyhere{
``A big move today'' and ``a big move yesterday'' are the two events whose
relationship measures volatility clustering. Making them precise --- a return
exceeding two standard deviations in absolute value --- is what turns an intuition
into something a generator can be scored against.
}

%% ─────────────────────────────────────────────────────────────────────────────
\section{The Probability Axioms}
\label{sec:axioms}

\situation{
You are asked to bet on a horse race. Someone quotes you odds implying a 60\%
chance for one horse and a 70\% chance for another, and there are five horses. You
do not need to know anything about horses to know these numbers are broken: they
already sum to more than certainty.
}

\problem{
Assigning numbers to events is only useful if the assignment is coherent. Three
rules are enough to rule out every incoherent assignment, and everything else in
probability follows from them.
}

\workedarith{
Fair die, each outcome probability $1/6$.

\medskip
\textbf{Non-negative:} every outcome has probability $1/6 > 0$. \checkmark

\textbf{Total is one:} $6 \times \tfrac16 = 1$. \checkmark

\textbf{Disjoint events add:} $A = \{2,4,6\}$ and $C = \{1,3\}$ share nothing, so
\[
  \Prob(A \cup C) = \Prob(\{1,2,3,4,6\}) = \tfrac56 ,
\]
and $\Prob(A) + \Prob(C) = \tfrac36 + \tfrac26 = \tfrac56$. \checkmark

\medskip
\textbf{Overlapping events do not simply add.} With $B = \{5,6\}$:
\[
  \Prob(A) + \Prob(B) = \tfrac36 + \tfrac26 = \tfrac56 ,
\]
but $A \cup B = \{2,4,5,6\}$ has probability $\tfrac46$. The excess $\tfrac16$ is
exactly $\Prob(A \cap B) = \Prob(\{6\})$, so
\[
  \Prob(A \cup B) = \tfrac36 + \tfrac26 - \tfrac16 = \tfrac46 . \quad\checkmark
\]

\medskip
\textbf{Complement.} $\Prob(A^{c}) = 1 - \Prob(A) = 1 - \tfrac36 = \tfrac36$, and
$A^{c} = \{1,3,5\}$ indeed has three outcomes. \checkmark
}

\formaldef{
A \textbf{probability measure} $\Prob$ assigns a number to each event subject to:
\begin{enumerate}
  \item $\Prob(A) \geq 0$ for every event $A$;
  \item $\Prob(\Omega) = 1$;
  \item if $A_1, A_2, \ldots$ are pairwise disjoint, then
        $\Prob\left(\bigcup_i A_i\right) = \sum_i \Prob(A_i)$.
\end{enumerate}
These are the \textbf{Kolmogorov axioms}. Immediate consequences:
$\Prob(A^{c}) = 1 - \Prob(A)$, $\Prob(\emptyset) = 0$, $\Prob(A) \leq 1$, and the
\textbf{inclusion--exclusion} rule
\[
  \Prob(A \cup B) = \Prob(A) + \Prob(B) - \Prob(A \cap B) .
\]
}

\analogybreak{
The axioms say what a coherent assignment looks like. They say nothing about
which assignment is \emph{correct} for any real situation --- that is a modelling
question, answered by data or judgement, not by mathematics. A Gaussian model of
returns satisfies every axiom perfectly and is still badly wrong about markets,
as \S\ref{sec:clt} shows.
}

\whyhere{
The axioms are what make the shuffled-real control a valid baseline. Shuffling
permutes outcomes without changing which values occur, so every probability
depending only on the collection of values is untouched --- which is precisely why
the control scores a Wasserstein distance of exactly $0.000$e$+00$ against real
data while failing every ordering-sensitive test.
}

%% ─────────────────────────────────────────────────────────────────────────────
\section{Conditional Probability}
\label{sec:conditional}

\situation{
A city has a 60\% chance of rain on any given day. But if it rained yesterday, the
chance today is 80\%. Knowing yesterday's outcome changes the probability you
assign to today's.
}

\problem{
Unconditional probabilities describe the long run and are silent about the
present. A risk model that knows only that big moves happen 5\% of the time, with
no way to use the fact that yesterday was violent, is discarding most of what the
data contains.
}

\workedarith{
Let $A$ = ``big market move today'', $B$ = ``big market move yesterday''.

MOEX, 20-year data: $\Prob(A) = 3.8\%$ while $\Prob(A \mid B) = 22.9\%$.

\medskip
\textbf{Ratio:} $22.9 / 3.8 = 6.0$. Knowing yesterday was violent multiplies
today's probability by roughly six.

\medskip
\textbf{Where the formula comes from.} Conditioning on $B$ restricts attention to
the days when $B$ happened, so $B$ becomes the new sample space. Of those days, we
ask what fraction also had $A$:
\[
  \Prob(A \mid B) = \frac{\Prob(A \cap B)}{\Prob(B)} .
\]

\textbf{On the die.} $A = \{2,4,6\}$, $B = \{5,6\}$:
\[
  \Prob(A \mid B) = \frac{\Prob(\{6\})}{\Prob(\{5,6\})}
  = \frac{1/6}{2/6} = \tfrac12 .
\]
Check by counting: told the roll exceeded four, it is 5 or 6, one of which is
even. One in two. \checkmark

\medskip
Note the asymmetry: $\Prob(B \mid A) = \frac{1/6}{3/6} = \tfrac13 \neq \tfrac12$.
}

\formaldef{
For $\Prob(B) > 0$,
\[
  \Prob(A \mid B) = \frac{\Prob(A \cap B)}{\Prob(B)} .
\]
Rearranging gives the \textbf{multiplication rule}
$\Prob(A \cap B) = \Prob(A \mid B)\,\Prob(B)$.

Events $A$ and $B$ are \textbf{independent} when $\Prob(A \mid B) = \Prob(A)$,
equivalently $\Prob(A \cap B) = \Prob(A)\Prob(B)$.
}

\analogybreak{
Conditioning is not causation. MOEX's clustering after February 2022 is driven by
a stream of news, not by the previous day's move causing the next. And the
relationship is about \emph{magnitude}, not direction: knowing yesterday was
violent says a great deal about how large today's move will be and essentially
nothing about its sign. That asymmetry is the whole content of volatility
clustering.
}

\whyhere{
The ratio $\Prob(A \mid B)/\Prob(A) = 6.0$ is a direct measurement of volatility
clustering, and reproducing it is what separates a generator that has learned
market dynamics from one that has only learned the histogram. The shuffled control
destroys exactly this structure: after permutation
$\Prob(A \mid B) \approx \Prob(A)$, the ratio collapses to 1, and the control fails
every temporal metric while acing every fidelity one.
}

%% ─────────────────────────────────────────────────────────────────────────────
\section{Bayes' Theorem}
\label{sec:bayes}

\situation{
A test for a rare disease is 99\% accurate. You test positive. The intuitive
reading --- ``99\% chance I have it'' --- is wrong, and not slightly. If the disease
affects one person in a thousand, the true figure is closer to 9\%.
}

\problem{
People reliably confuse $\Prob(\text{positive} \mid \text{ill})$ with
$\Prob(\text{ill} \mid \text{positive})$. The previous section already showed
these differ; Bayes' theorem says exactly how, and the gap is largest precisely
when the underlying event is rare --- which is the regime every tail-risk question
lives in.
}

\workedarith{
Take 100{,}000 people, one in a thousand ill, test 99\% accurate in both
directions.

\medskip
\begin{tabular}{@{}lrrr@{}}
\toprule
& Ill & Healthy & Total \\
\midrule
Test positive & $99$ & $999$   & $1{,}098$ \\
Test negative & $1$  & $98{,}901$ & $98{,}902$ \\
\midrule
Total         & $100$ & $99{,}900$ & $100{,}000$ \\
\bottomrule
\end{tabular}

\medskip
Ill: $100$ people, of whom $99\%$ test positive $\to 99$.

Healthy: $99{,}900$ people, of whom $1\%$ test positive $\to 999$ false alarms.

\medskip
\textbf{Given a positive test:}
\[
  \Prob(\text{ill} \mid +) = \frac{99}{99 + 999} = \frac{99}{1098} = 0.0902 = 9.0\% .
\]

The false alarms outnumber the true positives ten to one, because there are a
thousand times more healthy people to generate them. \checkmark

\medskip
\textbf{Via the formula:}
\[
  \Prob(\text{ill} \mid +) = \frac{\Prob(+ \mid \text{ill})\Prob(\text{ill})}{\Prob(+)}
  = \frac{0.99 \times 0.001}{0.01098} = 0.0902 . \quad\checkmark
\]
}

\formaldef{
\textbf{Bayes' theorem.} For events with $\Prob(B) > 0$,
\[
  \Prob(A \mid B) = \frac{\Prob(B \mid A)\,\Prob(A)}{\Prob(B)} .
\]
The denominator expands by the \textbf{law of total probability}:
\[
  \Prob(B) = \Prob(B \mid A)\Prob(A) + \Prob(B \mid A^{c})\Prob(A^{c}) .
\]
$\Prob(A)$ is the \textbf{prior}, $\Prob(A \mid B)$ the \textbf{posterior}, and
$\Prob(B \mid A)$ the \textbf{likelihood}.
}

\analogybreak{
Bayes' theorem is an identity, not a philosophy. It tells you how to update given
a prior; it does not tell you what prior to hold, and where the prior is
contested the theorem settles nothing. The disease example is clean only because
the base rate of one in a thousand was handed over as fact.
}

\whyhere{
Base rates dominate rare-event inference, and every tail question in this thesis
is a rare-event question. A test that flags a bad generator with 99\% accuracy
still produces mostly false alarms if bad generators are rare --- which is the same
arithmetic that makes multiple testing across 19 metrics
(\S\ref{sec:multipletesting}) so treacherous.
}

%% ─────────────────────────────────────────────────────────────────────────────
\section{Random Variables}
\label{sec:randomvars}

\situation{
Roll two dice and add them. The outcome of the experiment is a pair --- say
$(3,4)$ --- but what you care about is the single number 7. You have attached a
number to each outcome.
}

\problem{
Events are collections of outcomes, and outcomes are often not numbers: heads or
tails, a pair of dice, a market state. To do arithmetic --- averages, spreads,
distances --- we need numbers, and we need a disciplined way of attaching them.
}

\workedarith{
Two dice, 36 equally likely outcomes. Let $X$ be the sum.

\medskip
\begin{tabular}{@{}lrrrrrrrrrrr@{}}
\toprule
$X$ & 2 & 3 & 4 & 5 & 6 & 7 & 8 & 9 & 10 & 11 & 12 \\
\midrule
Ways & 1 & 2 & 3 & 4 & 5 & 6 & 5 & 4 & 3 & 2 & 1 \\
\bottomrule
\end{tabular}

\medskip
The counts total $1+2+3+4+5+6+5+4+3+2+1 = 36$. \checkmark

$X = 7$ arises six ways --- $(1,6), (2,5), (3,4), (4,3), (5,2), (6,1)$ --- so
$\Prob(X = 7) = 6/36 = 1/6$, while $\Prob(X = 2) = 1/36$.

\medskip
$X$ is a \emph{function}: it takes an outcome and returns a number, exactly the
machine of \S\ref{sec:functions}. What makes it \emph{random} is that its input is
the result of a chance experiment, not the function itself.
}

\formaldef{
A \textbf{random variable} is a function $X : \Omega \to \R$ assigning a real
number to each outcome. It is \textbf{discrete} if its range is finite or
countable, \textbf{continuous} if it takes values across an interval.

We write $\Prob(X = x)$ for the probability of the event
$\{\omega \in \Omega : X(\omega) = x\}$, and $\Prob(X \leq x)$ similarly. Capital
letters denote the random variable, lower case a particular value it might take.
}

\analogybreak{
``Random variable'' is a poor name twice over: it is not a variable, and it is not
random. It is an ordinary, fully determined function. All the randomness lives in
which outcome the experiment produces --- the function merely translates that
outcome into a number.
}

\whyhere{
The daily log return $r_t$ is a random variable. So is every statistic computed
from it: the Hill estimator, the sample kurtosis, the discriminative AUC. Treating
these as random variables rather than fixed numbers is what makes it meaningful to
ask about their sampling variability --- the subject of Chapter~\ref{ch:statistics}.
}

%% ─────────────────────────────────────────────────────────────────────────────
\section{PMF, PDF and CDF}
\label{sec:distributions}

\situation{
Two ways to describe rainfall. You could say how likely each amount is --- ``2\,mm
is common, 40\,mm is rare''. Or you could say how likely it is to be \emph{at
most} a given amount --- ``90\% of days see 10\,mm or less''. Both describe the same
weather.
}

\problem{
For discrete quantities, listing the probability of each value works. For
continuous ones it collapses: the probability of a return being exactly
$0.0123456\ldots$ is zero, and that is true of every single value, so the list
carries no information. Continuous quantities need a different device.
}

\workedarith{
\textbf{Discrete: the PMF.} For the two-dice sum,
$\Prob(X=2) = 1/36$, $\Prob(X=7) = 6/36$, and the values sum to 1.

\medskip
\textbf{Discrete: the CDF.} Accumulate:
\[
  F(4) = \Prob(X \leq 4) = \tfrac{1}{36} + \tfrac{2}{36} + \tfrac{3}{36} = \tfrac{6}{36} = \tfrac16 .
\]
And $F(12) = 1$, since every outcome is at most 12.

\medskip
\textbf{Continuous: why point probabilities vanish.} Pick a real number uniformly
from $[0,1]$. The chance it lands in $[0.3, 0.4]$ is $0.1$; in $[0.30, 0.31]$ is
$0.01$; in an interval of width $w$ is $w$. Shrink $w$ to zero and the probability
goes to zero. Any single point has probability exactly zero, yet some point
certainly occurs.

\medskip
\textbf{Continuous: the density.} What survives is probability \emph{per unit
width}. For the uniform on $[0,1]$ the density is $f(x) = 1$ throughout, and
\[
  \Prob(0.3 \leq X \leq 0.4) = \int_{0.3}^{0.4} 1\,dx = 0.4 - 0.3 = 0.1 . \quad\checkmark
\]
The density is a rate; the probability is its integral (\S\ref{sec:integrals}).
This is why a density may exceed 1 --- a uniform on $[0, 0.5]$ has $f(x) = 2$ ---
without anything being wrong: it is a rate, not a probability.
}

\formaldef{
For a discrete $X$, the \textbf{probability mass function} is
$p(x) = \Prob(X = x)$, with $p(x) \geq 0$ and $\sum_x p(x) = 1$.

For a continuous $X$, the \textbf{probability density function} $f$ satisfies
\[
  \Prob(a \leq X \leq b) = \int_a^b f(x)\,dx , \qquad
  f(x) \geq 0 , \qquad \int_{-\infty}^{\infty} f(x)\,dx = 1 .
\]

In both cases the \textbf{cumulative distribution function} is
\[
  F(x) = \Prob(X \leq x) ,
\]
non-decreasing, with $F(-\infty) = 0$ and $F(+\infty) = 1$. For continuous $X$ the
two are linked by the Fundamental Theorem of Calculus (\S\ref{sec:ftc}):
$F' = f$, and $F$ is the integral of $f$.
}

\analogybreak{
A density is not a probability and cannot be compared across different
quantities without care: change the units of $x$ and $f$ changes too, since it is
measured per unit of $x$. The CDF has no such problem --- it is a probability, and
dimensionless. This is one reason distances between distributions in this thesis
are built from CDFs and quantiles rather than densities.
}

\whyhere{
The Kolmogorov--Smirnov statistic is the largest vertical gap between two CDFs.
Value at Risk is an inverse CDF --- the return below which only 5\% of days fall.
And the histograms in the stylised-facts figure are estimates of $f$; the
log-scaled panel (\S\ref{sec:logscale}) is the one that makes the tail of $f$
legible.
}

%% ─────────────────────────────────────────────────────────────────────────────
\section{Expectation}
\label{sec:expectation}

\situation{
A fairground game costs \$2. You win \$10 with probability $1/8$ and nothing
otherwise. Play once and you either win or lose. Play ten thousand times and the
outcome is no longer in doubt: you will have lost money, and by a predictable
amount.
}

\problem{
Single outcomes are unpredictable; long-run averages are not. We need the number
that long-run averages converge to, computed from the distribution rather than
from data.
}

\workedarith{
\textbf{The game.} Win \$10 with probability $1/8$, \$0 with probability $7/8$:
\[
  \E[\text{winnings}] = 10 \times \tfrac18 + 0 \times \tfrac78 = 1.25 .
\]
Cost \$2, so the expected profit is $1.25 - 2 = -0.75$ per play. Over 10{,}000
plays, expect to lose about \$7{,}500.

\medskip
\textbf{One die.}
\[
  \E[X] = \tfrac16(1+2+3+4+5+6) = \tfrac{21}{6} = 3.5 .
\]
Note that 3.5 is not a possible outcome. The expectation is a balance point, not
a prediction.

\medskip
\textbf{Two dice.} By the same weighting,
\[
  \E[X] = \tfrac{1}{36}\bigl(2 \cdot 1 + 3 \cdot 2 + \cdots + 12 \cdot 1\bigr)
  = \tfrac{252}{36} = 7 .
\]
Which is $3.5 + 3.5$ --- expectation adds. \checkmark

\medskip
\textbf{Linearity, checked.} Let $Y = 2X + 1$ for one die. Directly:
$Y$ takes values $3,5,7,9,11,13$, averaging $\tfrac{48}{6} = 8$. By the rule:
$2\E[X] + 1 = 2(3.5) + 1 = 8$. \checkmark
}

\formaldef{
The \textbf{expectation} of a random variable is
\[
  \E[X] = \sum_x x\,p(x) \quad\text{(discrete)}, \qquad
  \E[X] = \int_{-\infty}^{\infty} x f(x)\,dx \quad\text{(continuous)} .
\]
It is \textbf{linear}: for constants $a, b$ and any $X, Y$,
\[
  \E[aX + bY] = a\E[X] + b\E[Y] ,
\]
which holds whether or not $X$ and $Y$ are independent.

For a function of $X$, $\E[g(X)] = \sum_x g(x)p(x)$. In general
$\E[g(X)] \neq g(\E[X])$.
}

\analogybreak{
Expectation requires the defining sum or integral to converge, and for
heavy-tailed distributions it may not. If $\Prob(|X| > x) \sim x^{-\alpha}$ with
$\alpha \leq 1$, then $\E[X]$ simply does not exist --- not ``is large'' or ``is
hard to estimate'', but is undefined. Our markets have $\hat\alpha$ between 2.54
and 3.21, so the mean and variance are fine, but this is a property to be checked
rather than assumed, and \S\ref{sec:clt} shows which moments do fail.
}

\whyhere{
Every loss function minimised in this thesis is an expectation estimated by a
batch average. The WGAN critic loss $\E[D(x_{\text{real}})] - \E[D(x_{\text{fake}})]$
and the gradient penalty $\E[(\norm{\nabla D}_2 - 1)^2]$ are both expectations,
approximated by averaging over a batch of 64 or 128 windows.
}

%% ─────────────────────────────────────────────────────────────────────────────
\section{Variance and Higher Moments}
\label{sec:variance}

\situation{
Two investments both average 5\% a year. One returns between 4\% and 6\% every
year; the other alternates $+40\%$ and $-30\%$. The average is a wholly inadequate
description of the difference.
}

\problem{
Location without spread is half a summary. And for financial returns even spread
is not enough: the shape of the extremes matters more than their typical size, so
we need a ladder of increasingly fine descriptions --- and a clear account of when
each rung exists.
}

\workedarith{
\textbf{One die.} $\E[X] = 3.5$. Deviations and their squares:

\medskip
\begin{tabular}{@{}lrrrrrr@{}}
\toprule
$x$ & 1 & 2 & 3 & 4 & 5 & 6 \\
$x - 3.5$ & $-2.5$ & $-1.5$ & $-0.5$ & $0.5$ & $1.5$ & $2.5$ \\
$(x-3.5)^2$ & $6.25$ & $2.25$ & $0.25$ & $0.25$ & $2.25$ & $6.25$ \\
\bottomrule
\end{tabular}

\medskip
\[
  \Var(X) = \tfrac16(6.25 + 2.25 + 0.25 + 0.25 + 2.25 + 6.25)
  = \tfrac{17.5}{6} = 2.9167 ,
\]
so the standard deviation is $\sqrt{2.9167} = 1.708$.

\medskip
\textbf{The shortcut.} $\E[X^2] = \tfrac16(1+4+9+16+25+36) = \tfrac{91}{6} = 15.1667$, and
\[
  \Var(X) = \E[X^2] - (\E[X])^2 = 15.1667 - 12.25 = 2.9167 . \quad\checkmark
\]

\medskip
\textbf{Why squares, not absolute values.} Squaring makes variance add for
independent variables, which absolute deviation does not. For two independent
dice, $\Var(\text{sum}) = 2.9167 + 2.9167 = 5.8333$.

\medskip
\textbf{Scaling.} $\Var(aX) = a^2\Var(X)$: doubling a return quadruples its
variance but only doubles its standard deviation. This is why volatility is
quoted as a standard deviation --- it shares the units of the returns themselves.
}

\formaldef{
The \textbf{variance} is
\[
  \Var(X) = \E\bigl[(X - \E[X])^2\bigr] = \E[X^2] - (\E[X])^2 ,
\]
and the \textbf{standard deviation} is $\sigma = \sqrt{\Var(X)}$.

The \textbf{$k$-th moment} is $\E[X^k]$. Standardised third and fourth moments
give
\[
  \text{Skewness} = \frac{\E[(X-\mu)^3]}{\sigma^3} , \qquad
  \text{Kurtosis} = \frac{\E[(X-\mu)^4]}{\sigma^4} ,
\]
with excess kurtosis defined as kurtosis minus 3, so that a normal distribution
scores 0.

Properties: $\Var(aX + b) = a^2\Var(X)$, and $\Var(X+Y) = \Var(X) + \Var(Y)$ when
$X$ and $Y$ are independent.
}

\analogybreak{
Higher moments exist only when their integrals converge. For a power-law tail
$\Prob(|X|>x) \sim x^{-\alpha}$, the $k$-th moment is finite only if $k < \alpha$.
The BRICS markets in this project have $\hat\alpha$ between $2.54$ and $3.21$:
\begin{itemize}
  \item $k=2$ (variance): finite everywhere, since $2 < 2.54$;
  \item $k=3$ (skewness): infinite in four of five markets;
  \item $k=4$ (kurtosis): \emph{infinite in all five}.
\end{itemize}
The population kurtosis of these return series does not exist. A sample kurtosis
can always be computed, but it is estimating nothing.
}

\whyhere{
This is why \texttt{kurtosis\_diff} and \texttt{skewness\_diff} were moved out of
the ranked metrics and into the descriptive-only group. Ranking models on an
estimate of a quantity that is infinite is not a strict test --- it is noise with a
decimal point, and Chapter~\ref{ch:heavytails} shows the estimate diverging with
sample size rather than converging.
}

%% ─────────────────────────────────────────────────────────────────────────────
\section{Joint Distributions and Independence}
\label{sec:joint}

\situation{
Height and weight, measured across a population. Knowing someone is tall changes
what you expect their weight to be. Height and shoe colour, by contrast, tell you
nothing about each other.
}

\problem{
Individual distributions cannot express a relationship. Two return series may have
identical histograms and behave completely differently in combination --- and the
whole question of temporal structure is a question about the joint behaviour of
$r_t$ and $r_{t-1}$.
}

\workedarith{
Two dice. Let $X$ be the first roll and $S$ the sum.

\medskip
\textbf{Are the two rolls independent?} $\Prob(\text{first} = 3) = 1/6$ and
$\Prob(\text{second} = 5) = 1/6$; the pair $(3,5)$ is one of 36 outcomes, so
\[
  \Prob(3, 5) = \tfrac{1}{36} = \tfrac16 \times \tfrac16 . \quad\checkmark
\]
The joint factorises, so they are independent.

\medskip
\textbf{Are $X$ and $S$ independent?} $\Prob(X = 1) = 1/6$ and
$\Prob(S = 12) = 1/36$. If independent, $\Prob(X=1, S=12)$ would be
$\tfrac16 \times \tfrac{1}{36} = \tfrac{1}{216}$. But $S = 12$ requires both dice
to be 6, so $X = 1$ and $S = 12$ cannot both happen:
\[
  \Prob(X=1, S=12) = 0 \neq \tfrac{1}{216} .
\]
Not independent. \checkmark Knowing the sum tells you about the first roll.

\medskip
\textbf{Conditional distribution.} Given $S = 4$, the possibilities are
$(1,3), (2,2), (3,1)$, so $X \in \{1,2,3\}$ each with probability $1/3$ --- a
completely different distribution from the unconditional uniform on $\{1,\ldots,6\}$.
}

\formaldef{
The \textbf{joint distribution} of $X$ and $Y$ gives $\Prob(X = x, Y = y)$, or in
the continuous case a joint density $f(x,y)$ with
\[
  \Prob\bigl((X,Y) \in A\bigr) = \iint_A f(x,y)\,dx\,dy .
\]
The \textbf{marginal} distribution of $X$ recovers $X$ alone by summing or
integrating out $Y$.

$X$ and $Y$ are \textbf{independent} if the joint factorises for all values:
\[
  \Prob(X = x, Y = y) = \Prob(X = x)\Prob(Y = y) ,
  \qquad\text{or}\qquad f(x,y) = f_X(x)f_Y(y) .
\]
Equivalently, the conditional distribution of $X$ given $Y$ does not depend on $Y$.
}

\analogybreak{
Marginals do not determine the joint. Two pairs of variables can have identical
individual distributions and utterly different dependence --- this is exactly the
gap a copula is designed to describe. It is also the precise sense in which
matching a histogram is not matching a process, and the reason the evaluation
framework in this thesis needs a second family of metrics at all.
}

\whyhere{
The shuffled-real control has \emph{identical} marginals to the real data by
construction and destroyed joint structure. That single fact is why an evaluation
built only on marginal distances ranked it first: under the old unweighted
19-metric composite it scored an average rank of $1.24$ against $2.47$ for a
genuine generator. The permutation-invariant metrics could not see the difference,
because there is no difference for them to see.
}

%% ─────────────────────────────────────────────────────────────────────────────
\section{Covariance and Correlation}
\label{sec:covariance}

\situation{
Ice cream sales and drowning deaths rise together over the year. Neither causes
the other --- both follow temperature --- but the co-movement is real and
measurable.
}

\problem{
Independence is a yes/no property. What is usually wanted is a graded measure:
\emph{how strongly} do two quantities move together, on a scale that permits
comparison across quantities with different units?
}

\workedarith{
Four paired observations:

\medskip
\begin{tabular}{@{}lrrrr@{}}
\toprule
$x$ & 1 & 2 & 3 & 4 \\
$y$ & 2 & 4 & 5 & 9 \\
\bottomrule
\end{tabular}

\medskip
Means: $\bar{x} = 10/4 = 2.5$, $\bar{y} = 20/4 = 5$.

\medskip
\begin{tabular}{@{}lrrrr@{}}
\toprule
$x - \bar{x}$ & $-1.5$ & $-0.5$ & $0.5$ & $1.5$ \\
$y - \bar{y}$ & $-3$   & $-1$   & $0$   & $4$ \\
product       & $4.5$  & $0.5$  & $0$   & $6$ \\
\bottomrule
\end{tabular}

\medskip
\[
  \Cov(x,y) = \tfrac14(4.5 + 0.5 + 0 + 6) = \tfrac{11}{4} = 2.75 .
\]
Positive: the two rise together.

\medskip
\textbf{To correlation.} Spreads:
$\Var(x) = \tfrac14(2.25+0.25+0.25+2.25) = 1.25$, so $\sigma_x = 1.118$;
$\Var(y) = \tfrac14(9+1+0+16) = 6.5$, so $\sigma_y = 2.550$.
\[
  \rho = \frac{2.75}{1.118 \times 2.550} = \frac{2.75}{2.851} = 0.965 .
\]
Close to the maximum of 1 --- a strong, nearly linear relationship.

\medskip
\textbf{Units cancel.} Measure $y$ in different units, say doubled: covariance
doubles to $5.5$, but $\sigma_y$ doubles too, so $\rho$ is unchanged. \checkmark
}

\formaldef{
\[
  \Cov(X,Y) = \E\bigl[(X - \E X)(Y - \E Y)\bigr] = \E[XY] - \E[X]\E[Y] ,
\]
\[
  \rho(X,Y) = \frac{\Cov(X,Y)}{\sigma_X \sigma_Y} \in [-1, 1] .
\]
$\rho = \pm 1$ exactly when $Y$ is an increasing or decreasing linear function of
$X$. If $X$ and $Y$ are independent then $\Cov(X,Y) = 0$.

The \textbf{autocorrelation} of a series at lag $k$ is the correlation of the
series with itself shifted by $k$:
\[
  \rho(k) = \frac{\Cov(r_t, r_{t-k})}{\Var(r_t)} .
\]
}

\analogybreak{
Zero correlation does not imply independence. Let $X$ be symmetric about zero and
$Y = X^2$; then $\Cov(X,Y) = 0$ while $Y$ is completely determined by $X$.
Correlation detects \emph{linear} association only.

This is exactly the situation for financial returns. The autocorrelation of $r_t$
is near zero --- returns are close to linearly unpredictable --- while the
autocorrelation of $|r_t|$ and $r_t^2$ is strongly positive and decays slowly.
Returns are uncorrelated and emphatically not independent, which is the single
most important structural fact in this thesis.
}

\whyhere{
This is why the evaluation computes three separate autocorrelation metrics ---
\texttt{acf\_returns\_mae} on $r_t$, \texttt{acf\_absolute\_mae} on $|r_t|$, and
\texttt{acf\_squared\_mae} on $r_t^2$. A generator that matches only the first has
reproduced the absence of linear predictability while missing volatility
clustering entirely, and only the second and third catch it.
}

%% ─────────────────────────────────────────────────────────────────────────────
\section{The Law of Large Numbers}
\label{sec:lln}

\situation{
Flip a fair coin ten times and you might see seven heads --- unremarkable. Flip it
ten thousand times and 7{,}000 heads would be extraordinary. The proportion
settles toward one half, and it settles more firmly the longer you go.
}

\problem{
Expectation was defined from a distribution, and we asserted that long-run
averages converge to it. That assertion needs to be a theorem, with stated
conditions --- because the conditions are exactly what financial data threatens.
}

\workedarith{
Rolling a fair die, $\E[X] = 3.5$. The standard deviation of the average of $n$
rolls is $\sigma/\sqrt{n}$ with $\sigma = 1.708$:

\medskip
\begin{tabular}{@{}lrr@{}}
\toprule
$n$ & sd of the average & typical spread \\
\midrule
$1$      & $1.708$  & $\pm 1.71$ \\
$100$    & $0.171$  & $\pm 0.17$ \\
$10{,}000$ & $0.017$ & $\pm 0.02$ \\
\bottomrule
\end{tabular}

\medskip
Each hundredfold increase in $n$ shrinks the spread tenfold, since
$\sqrt{100} = 10$. Convergence is real but slow: to halve the uncertainty you must
quadruple the data.

\medskip
\textbf{The cost of that slowness.} This project's per-market test set holds about
496 observations; pooling all five markets gives about 2{,}480. The pooled
standard error is $\sqrt{2480/496} = \sqrt{5} = 2.24$ times smaller --- which is why
the downstream utility metrics are computed on the pooled set rather than
per-market.
}

\formaldef{
\textbf{Law of Large Numbers.} If $X_1, X_2, \ldots$ are independent with common
distribution and $\E|X| < \infty$, then the sample mean
\[
  \bar{X}_n = \frac{1}{n}\sum_{i=1}^n X_i \longrightarrow \E[X]
\]
as $n \to \infty$. When the variance $\sigma^2$ is also finite,
$\Var(\bar{X}_n) = \sigma^2/n$, so the standard error is $\sigma/\sqrt{n}$.
}

\analogybreak{
The law requires the mean to exist. It says nothing about how fast convergence
happens for any particular distribution, and for heavy-tailed data the answer can
be ``far more slowly than $1/\sqrt{n}$ suggests'', because a single extreme
observation can move the running average substantially even at large $n$. The
guarantee is asymptotic; the sample you actually hold is finite.
}

\whyhere{
Every metric in this thesis is an average over a finite sample and therefore
carries a standard error. The discriminative AUC null is not exactly $0.5$ but
$0.506 \pm 0.084$ --- measured across fifteen real-versus-real splits --- and that
spread is precisely the $\sigma/\sqrt{n}$ of this section. An observed AUC must be
judged against that spread, not against the theoretical $0.5$.
}

%% ─────────────────────────────────────────────────────────────────────────────
\section{The Central Limit Theorem --- and Where It Fails}
\label{sec:clt}

\situation{
Roll a die 100 times and record the average. Repeat the whole experiment 10{,}000
times and plot the 10{,}000 averages. A bell curve appears --- even though a single
die is flat, with all six faces equally likely. The bell was not in the die; it
came from the averaging.
}

\problem{
The Central Limit Theorem is the reason the normal distribution is everywhere, and
it is routinely over-read as licensing Gaussian models for anything. Applied to
daily market returns it makes predictions that are wrong by orders of magnitude,
and understanding exactly which of its conditions fails --- and which does not --- is
the intellectual core of this thesis.
}

\workedarith{
\textbf{What the Gaussian predicts.} For a normal distribution,
\[
  \Prob(|Z| > 5) = 2\Phi(-5) \approx 5.73 \times 10^{-7} .
\]
At 252 trading days per year, the expected waiting time for one such day is
\[
  \frac{1}{252 \times 5.73 \times 10^{-7}} \approx 6{,}922 \text{ years} .
\]

\medskip
\textbf{What the data shows.} A single BRICS market contributes roughly 6{,}201
market-days --- about 25 years. The Gaussian expects
$6201 \times 5.73\times10^{-7} \approx 0.0036$ five-sigma days over that span:
call it one day every three centuries. Eleven were observed.

The discrepancy is a factor of roughly $11 / 0.0036 \approx 3{,}000$. This is not a
model that is slightly miscalibrated in the tails; it is a model whose tail
predictions are meaningless.

\medskip
\textbf{Which condition failed?} The CLT needs finite variance. Our markets have
$\hat\alpha \in [2.54, 3.21] > 2$, so the variance \emph{is} finite and the CLT
\emph{does} apply to averages of returns. The theorem is not violated.

What fails is the inference people draw from it. The CLT describes the
distribution of the \emph{average} of many returns. It says nothing whatever about
the distribution of a \emph{single} return, and a five-sigma day is a single
observation, not an average. Modelling individual daily returns as Gaussian is not
an application of the CLT --- it is an unrelated assumption that happens to be
false.

\medskip
\textbf{And the fourth moment.} Even where the CLT applies, it applies to the
mean. The analogous result for the sample variance requires a finite fourth
moment, and with $\alpha < 4$ everywhere in our data that moment is infinite. So
the sample kurtosis has no limiting normal distribution, no standard error, and no
convergence --- Chapter~\ref{ch:heavytails} shows it climbing from $2.53$ at
$n = 250$ to $11.84$ at $n = 2{,}000$ on BOVESPA while the Hill estimator, which
does not require finite moments, settles from $4.83$ to $3.11$.
}

\formaldef{
\textbf{Central Limit Theorem.} Let $X_1, X_2, \ldots$ be independent with common
distribution, mean $\mu$ and \emph{finite} variance $\sigma^2$. Then
\[
  \frac{\bar{X}_n - \mu}{\sigma/\sqrt{n}} \longrightarrow \N(0,1)
\]
in distribution as $n \to \infty$.

Three conditions carry the weight: independence, identical distribution, and
finite variance. When the variance is infinite ($\alpha < 2$), averages converge
instead to a \textbf{stable} distribution with heavy tails of their own, and the
normal approximation never appears at any sample size.
}

\analogybreak{
The CLT is about sums and averages, not about individual observations, and this is
the confusion it most often causes. It also assumes independence, which market
returns violate through volatility clustering --- versions of the theorem exist for
dependent sequences, but they converge more slowly and require their own
conditions.

So there are two distinct failures to keep apart. Applying the CLT to a single
day's return is a misapplication of a correct theorem. Computing a sample kurtosis
and treating it as an estimate is applying a theorem whose conditions genuinely do
not hold.
}

\whyhere{
This section is the motivation for the entire thesis. If daily returns were
Gaussian, a closed-form risk model would suffice and no generative model would be
needed. Because they are not --- eleven five-sigma days where the model expects
0.0036 --- we need generators that reproduce heavy tails and volatility clustering,
and an evaluation framework honest enough to check whether they actually do. It is
also why the tail index $\hat\alpha$, which needs no finite moments, is a ranked
metric while kurtosis is not.
}
